\section{Auth0 Acquisition：共享 Primitive，不急于合并所有边界}

McKinnon 把 Okta 与 Auth0 描述为互补：前者长期面向 workforce/IT，后者以 developer-first customer identity 著称。两者共享 directory、authentication、authorization、protocol、security 和 scale primitive，却服务不同 buyer、end user、SDK、customization 与 release expectation。并购架构的第一原则不是尽快画成一张 org chart，而是保护 customer contract。

\subsection{Workforce Identity 与 Customer Identity}

Workforce identity 关注员工、contractor、IT administration、enterprise app 和 lifecycle；customer identity 关注产品终端用户、developer integration、custom UX、social login、application-specific profile 与互联网流量。共同平台可以复用 protocol、threat intelligence 和 reliability practice，但 runtime migration、tenant model 与 developer experience 需要独立验证。

\lecturefigure{14-two-platform-acquisition.png}{并购应尊重 workforce/customer identity 的不同用户、工作流和基础设施边界。}{本地字幕 30:00--36:20；Okta/Auth0 2021 官方并购资料；概念重绘。}

\begin{knowledgebox}{读图：Shared primitive 不等于 shared runtime}
两类产品都需要 directory、OAuth/OIDC、MFA、security 和 scale；但 workforce buyer 可能要求 IT policy、employee lifecycle 与 enterprise integration，customer identity 则强调 SDK、品牌化体验、极端峰值与 developer control。并购后应先统一 threat intelligence、support escalation 和 strategic interface，再用 compatibility、SLO 和 migration evidence 决定是否合并 runtime。
\end{knowledgebox}

\teachervoice{讲者回顾并购时认为，投资者对 M\&A 的怀疑比自己预期更强；他也承认在低利率和高速增长环境下，部分业务整合可能推进过快。这个反思提醒系统团队：短期增长会掩盖 integration cost，不能把宏观顺风误当成架构已经验证。}

\begin{warningbox}{Platform consolidation 的三类隐藏成本}
统一销售包装可能先于 entitlement 和 billing 一致；统一 console 可能掩盖不同 tenant/security model；统一 backend 则会扩大 blast radius、迁移状态并改变客户合规边界。每一步都需要 rollback、data migration validation、customer communication 和 parallel run，而不是以“one company”替代技术验收。
\end{warningbox}

\subsection{本章小结}

并购整合是 product contract、runtime state 与组织激励的共同迁移。共享 identity primitive 可以产生协同，但 preserve boundary 往往比立即统一更安全。

